\documentstyle[% 


righttag% 


,11pt% 


,amssymb% 


,russian%               remove in Eng. 


,izvru%                 Set izven% in Eng. 


% ,izvru-ks             for brief communications 


]{amsart} 





 


%       Math definitions 


\newcommand{\const}{\operatorname{const}}


\newcommand{\bsy}{\boldsymbol}


\newcommand{\tts}[1]{}


 


%\hoffset=-15mm                  For Eng. 


 


\setcounter{page}{21} 


\god{1997} 


\nomer{9~(424)} %                Remove in Eng. 


%\nomer{Vol.\,41, No.\,9}        Set in Eng. 


% \str{75-81}                    Set in Eng. 


\udk{514.76} 


\author{\it П.Е.\,Марков} 


\title{Общие аналитические и\\


бесконечно малые деформации погружений. I} 


 


\date{} 


% \pagestyle{myheadings}         Set in Eng. 


\pagestyle{plain} %              Remove in Eng. 


\begin{document} 


% \thispagestyle{plain}          Set in Eng. 


\maketitle 


 


\section*{Введение}





Работа посвящена бесконечно малым (б.\,м.) деформациям конечного порядка


и аналитическим деформациям погружений гладкого $n$-мерного


многообразия $X$ в плоское $m$-мерное, $1\le n\le m$, пространство


$\Pi$ произвольной сигнатуры без каких-либо ограничений на вид


возникающих вариаций метрики и других сечений тензорных расслоений.


Такие деформации мы называем {\it общими} деформациями. Теория общих


деформаций представляет собою основу всех ``специальных'' теорий


деформаций: аналитических и б.\,м. изгибаний


(\cite{Efi}--\cite{Iva2}),


конформных деформаций \cite{Car1}, \cite{DoC},


деформаций с сохранением грассманова


образа \cite{Fom}, ареальных


деформаций \cite{Bes} и


т.\,д. (см. также \cite{Yano}, \cite{Chen}).


Несмотря на большую общность, эта теория оказывается богатой


содержанием. Так, для общей аналитической (б.\,м.) деформации погружения


$X\to\Pi$ можно построить систему полей вращений, вывести систему


основных уравнений (аналог основных уравнений теории поверхностей) и


доказать аналог основной теоремы теории поверхностей. Точные


формулировки этих результатов требуют введения некоторых обозначений и


определений и будут приведены позже. Основные обозначения,


терминология и предварительные результаты приводятся в \S\,1. В этом же


параграфе выводятся уравнения Гаусса, Петерсона-Кодацци и Риччи для


$C^2$-погружений.





\S\,2 посвящен системе полей вращений для общей аналитической (б.\,м.)


деформации погружения $X\to\Pi$. Эта система является одним из


фундаментальных понятий теории б.\,м. и аналитических изгибаний


поверхностей. Для б.\,м. изгибаний первого порядка двумерной поверхности


в $E^3$ поле вращений было известно еще Г.\,Дарбу


(\cite{Dar}, гл.\,3, с.\,48--72). В случае б.\,м. изгибаний конечного


порядка и аналитических изгибаний двумерной поверхности в $E^3$


система полей вращений рассматривалась в статье Н.В.\,Ефимова


\cite{Efi}. Для б.\,м. изгибаний первого порядка $n$-мерной


поверхности в $\Pi$ поле вращений было введено в работе


\cite{Mark1} как поле бивекторов в $\Pi$. Для б.\,м. изгибаний


произвольного порядка (аналитических изгибаний) погружения $X\to\Pi$


система полей вращений была введена в работе \cite{Mark2}. Здесь


мы докажем существование и единственность системы полей вращений для


общей б.\,м. деформации произвольного порядка


погружения $X\to\Pi$ и установим некоторые свойства


этой системы. В частности, будут выведены формулы, выражающие


дифференциалы полей вращений через вариации форм локального корепера,


связности, погружения и кручения, обобщающие известные формулы для


производных поля вращений двумерной поверхности в $E^3$


(\cite{Efi}, \cite{Vek}, гл.\,5, \S\,6, с.\,337).





В \S\,3 мы рассмотрим систему уравнений, решения которой определяют


б.\,м. деформации конечного порядка погружения $X\to\Pi$. Эта система


может быть получена формальным варьированием основных уравнений теории


поверхностей и потому называется основной системой теории


б.\,м. деформаций погружений. Проблема заключается в доказательстве


того, что для всякого решения этой системы найдется б.\,м. деформация


погружения, при которой вариации форм локального корепера, связности,


погружения и кручения совпадут с формами этого решения (аналог


основной теоремы теории поверхностей). В случае б.\,м. изгибаний первого


порядка двумерной поверхности в $E^3$ эта теорема была известна еще


Г.\,Дарбу (\cite{Dar}, гл.\,2, с.\,18--47). Для б.\,м. изгибаний первого


порядка $n$-мерных поверхностей в $m$-мерных пространствах постоянной


кривизны она доказана в работе \cite{Mark3}. Для б.\,м. изгибаний


высших порядков погружения $X\to\Pi$ --- в работе \cite{Mark2}.


Для б.\,м. деформаций первого порядка с заданной вариацией метрики


$n$-мерной поверхности в $m$-мерном евклидовом пространстве --- в


работе \cite{Mark4}. Здесь мы докажем эту теорему для общей


б.\,м. деформации произвольного конечного порядка погружения $X\to\Pi$.


В этом же параграфе устанавливается вид вариаций форм локального


корепера и связности, для которых соответствующая б.\,м. деформация


является б.\,м. изгибанием.





В \S\,4 проводится исследование основной системы уравнений в


зависимости от типового числа погружения. Типовое число было введено


К.Б.\,Аллендорфером \cite{All} в связи с двумя


фундаментальными результатами. Первый из них утверждает жесткость


погружений с типовым числом $\ge3$, второй --- что для погружений с


типовым числом $\ge4$ уравнения Петерсона-Кодацци и Риччи являются


следствиями уравнений Гаусса. Данное Аллендорфером определение


типового числа довольно громоздко и многими авторами подвергалось


обработке (см. \cite{Yan}--\cite{Kob}).


Мы докажем два


аналога второго из результатов Аллендорфера для рассмотренной в \S\,3


основной системы уравнений. Определение типового числа мы даем по


книге \cite{Kob} (Прим.\,17, c.\,317). Схему доказательств


заимствуем из статьи \cite{Chern}, где дается упрощенное


доказательство теоремы Аллендорфера.





\section{Предварительные результаты}





{\bf 1.1.} Всюду под $X$ будем понимать связное $n$-мерное, $n\ge 1$,


хаусдорфово


ориентируемое $C^{\infty}$-многообразие, удовлетворяющее второй


аксиоме счетности. Всякое расслоение с базой $X$ и тотальным


пространством $P(X)$ будем обозначать символом его тотального


пространства $P(X)$. Слой над точкой $x\in X$ будем обозначать через


$P_x(X)$. Через $T(X)$ обозначаем касательное к $X$ расслоение, через


$T^*(X)$ --- кокасательное к $X$ расслоение.


Рассмотрим путь


$\gamma:(-\varepsilon,\varepsilon)\to X$, $\varepsilon>0$,


$\gamma(0)=x\in X$, класса $C^s$, $s\ge 1$.


Обозначим через $C^s(x)$ множество всех функций класса $C^s$ в


некоторой окрестности точки $x$. Отображение $v:C^s(x)\to{\bold R}$,


определенное формулой


$v(f)=\left.\dfrac{d^sf(\gamma(t))}{dt^s}\right|_{t=0}$, называется


$s$-касательным вектором к пути $\gamma$ в точке $x$. Линейная


оболочка $T{}_x^{(s)}(X)$, натянутая на множество $s$-касательных


векторов к всевозможным $C^s$-путям, проходящим через точку $x$,


является векторным пространством размерности


$\sum\limits_{k=1}^{s}\dfrac{(n+k-1)!}{(n-1)!k!}$ и называется $s$-касательным


пространством к $X$ в точке $x$. Векторное расслоение, слоем над


каждой точкой $x\in X$ которого является $T{}_x^{(s)}(X)$, называется


$s$-касательным расслоением многообразия $X$ (этот термин мы


заимствовали из \cite{Vish}, гл.\,7, c.\,153). При $s=1$ получаем


касательное расслоение. Для всякого $C^r$-расслоения $P(X)$, $1\le s\le r$,


через $P^{(s)}(X)$ мы обозначаем $C^{r-s}$-расслоение, слоем над каждой точкой


$x\in X$ которого является $s$-касательное расслоение


$T^{(s)}(P_x(X))$ к слою $P_x(X)$. Расслоение $P^{(s)}(X)$ мы называем


$s$-производным расслоением расслоения $P(X)$.





Через ${\cal E}^r(X,P(X))$ будем обозначать множество всех


$C^r$-сечений расслоения $P(X)$. Для двух многообразий $X$ и $Y$ через


$C^r(X,Y)$ обозначаем множество всех отображений $X\to Y$ класса


$C^r$. Каждое отображение $f:X\to Y$ мы рассматриваем как сечение


$x \mapsto (x,f(x))$


тривиального расслоения $X\times Y$, так что


$C^r(X,Y)= {\cal E}^r(X,X\times Y)$.


\medskip





{\bf 1.2.} Пусть $\Im(X)$ --- расслоение локальных кореперов на $X$. Желая


подчеркнуть гладкость локального корепера, мы будем рассматривать его


как сечение этого расслоения, указывая область определения, в случае


необходимости, дополнительно. Пусть $G$ --- группа Ли, являющаяся


подгруппой полной линейной группы $\bold{GL(n)}$. О двух локальных


кореперах $\tau,\tau'\in{\cal E}^r(X,\Im(X))$ с областями


определения $U$, $U'$ соответственно будем говорить, что они


$G$-со\-гла\-со\-ва\-ны, если либо $U\cap U'=\emptyset$, либо в каждой точке


$x\in U\cap U'$ соответствующие базисы $\tau(x)=(\tau^i(x))_{i=1}^n$,


$\tau'(x)=(\tau'{}^k(x))_{k=1}^n$ пространства $T^*_x(X)$ связаны


равенством $\tau'{}^k(x)=p^k_i(x)\tau^i(x)$, где матрица


$P= (p^k_i)\in C^r(X,G)$. Множество всех попарно $G$-согласованных


локальных кореперов называется $G$-орбитой. Так как области


определения локальных кореперов покрывают $X$, то, хотя отношение


$G$-согласованности и не является отношением эквивалентности,


множество всех локальных $C^r$-кореперов на $X$ разбивается на


непересекающиеся $G$-орбиты, и каждая $G$-орбита однозначно


определяется заданием какого-либо ее корепера


(\cite{Zu}, гл.\,2, \S\,6, c.\,144). Всякая $G$-орбита


представляет собою локально-тривиальное расслоение со стандартным


слоем $G$ и структурной группой $G$.


\medskip





{\bf 1.3.} В качестве группы $G$ у нас будет выступать псевдоортогональная


группа, определяемая следующим образом. Зафиксируем упорядоченный


набор $\Delta=(\Delta^1,\dotsc,\Delta^n)$, где $\Delta^i=\pm 1$,


$i=1,\dotsc,n$. Будем обозначать также через $\Delta$ диагональную


$n\times n$-матрицу, на диагонали которой в $i$-й строке стоит


$\Delta^i$. Через $\bold{O(n,\bsy\Delta)}$ обозначим совокупность всех


$n\times n$-матриц $P$, удовлетворяющих условию


$\Delta P^T=P^{-1}\Delta$, где $P^T$ --- матрица, получаемая из $P$


транспонированием. Легко видеть, что $\bold{O(n,\bsy\Delta)}$ --- подгруппа


группы $\bold{GL(n)}$. Группу $\bold{O(n,\bsy\Delta)}$ будем называть


псевдоортогональной группой сигнатуры $\Delta$. Из ее определения


следует, что если $P\in{\bold{O(n,\bsy\Delta)}}$, то $\det P=\pm 1$.


Множество всех матриц $P\in{\bold{O(n,\bsy\Delta)}}$ с $\det P=+1$ образует


подгруппу группы $\bold{O(n,\bsy\Delta)}$, называемую специальной


псевдоортогональной группой и обозначаемую $\bold{SO(n,\bsy\Delta)}$. Тот


факт, что $\bold{SO(n,\bsy\Delta)}$ является группой Ли, выводится из


следующей леммы, распространяющей рациональную параметризацию


ортогональной группы по Кэли (\cite{Weyl}, гл.\,2, \S\,10) на группу


$\bold{SO(n,\bsy\Delta)}$, и используемую нами в дальнейшем. В формулировке


этой леммы ${\wedge(\bf n)}$ обозначает пространство


кососимметрических $n\times n$-матриц. Предполагается, что топология в


$\bold{SO(n,\bsy\Delta)}$ и ${\wedge(\bf n)}$ индуцируется евклидовой


нормой в пространстве всех $n\times n$-матриц (\cite{Gan},


гл.\,14, \S\,2, c.\,409).





\begin{lem}


Ненулевая матрица в пространстве ${\wedge(\bf n)}$ обладает


окрестностью $U$, которую отображение $\theta$, определенное формулой


$\theta(v)=(1+v\Delta)(1-v\Delta)^{-1}$, $v\in U$, аналитически


диффеоморфно отображает на окрестность $\theta(U)$ единичной матрицы


{\em 1} в группе $\bold{SO(n,\bsy\Delta)}$.


\end{lem}





\begin{pf} Обозначая через $U$ окрестность нулевой матрицы в


${\wedge(\bf n)}$, на которой $\det(1-v\Delta)\ne 0$, для всякой


матрицы $v\in U$ имеем


\begin{align*}


(\theta(v))^T &= (1+\Delta v)^{-1}(1-\Delta v)=((1-\Delta v)^{-1}+


(1-\Delta v)^{-1}\Delta v)^{-1}=\\


&=((1+\Delta v)(1-\Delta v)^{-1}+(1-\Delta v)(1-\Delta v)^{-1}-\\


&\quad -(1-\Delta v)^{-1}-(1-\Delta v)^{-1}(1-\Delta v)+


(1-\Delta v)^{-1})^{-1}=\\


&=((1+\Delta v)(1-\Delta v)^{-1})^{-1}=


\Delta(\theta(v))^{-1}\Delta.


\end{align*}


Следовательно, $\theta(v)\in{\bold{SO(n,\bsy\Delta)}}$. Легко видеть, что


обратное отображение $\theta^{-1}$ существует на $\theta(U)$ и имеет


вид $\theta^{-1}(P)=-(1-P)(1+P)^{-1}\Delta$, $P\in\theta(U)$. Так как


элементы матрицы $\theta(v)$ являются рациональными функциями элементов


матрицы $v$, а элементы матрицы $\theta^{-1}(P)$ --- рациональными


функциями элементов матрицы $P$, то $\theta$ --- аналитический


диффеоморфизм.


\end{pf}





Композицией отображения $\theta^{-1}$ со сдвигом в


$\bold{SO(n,\bsy\Delta)}$ мы можем ввести рациональную параметризацию в


окрестности любой матрицы $P\in{\bold{SO(n,\bsy\Delta)}}$ и, тем самым,


задать структуру аналитического многообразия на $\bold{SO(n,\bsy\Delta)}$,


согласованную с умножением в этой группе (\cite{Post},


Лек.\,1, c.\,17), а значит, превратить


$\bold{SO(n,\bsy\Delta)}$ в группу Ли.


\medskip





{\bf 1.4.} $\bold{SO(n,\bsy\Delta)}$-орбиту локального корепера


$\tau\in{\cal E}^r(X,\Im(X))$ будем обозначать через $\Re_{\Delta}(X)$.


Следующее предложение служит основой для построения тензорного анализа


над $\Re_{\Delta}(X)$.





\begin{lem}


Для всякого корепера


$\tau=(\tau^i)_{i=1}^n\in{\cal E}^r(X,\Re_{\Delta}(X))$,


$r\ge 1$, с областью определения


$U$ существует и притом единственная система $1$-форм


$\Phi^i_j\in{\cal E}^{r-1}(U,T^*(U))$, удовлетворяющая условиям


$$


d\tau^i=\tau^j\wedge \Phi^i_j,\quad \Delta^i\Phi^i_j+\Delta^j\Phi^j_i=0,


\quad i,j=1,\dotsc,n.


$$


\end{lem}





{\bf Доказательство} вытекает из теоремы 2.1.10 книги \cite{Wolf}


(гл.\,2, \S\,1, c.\,67--68) заменой обозначений. Формы $\Phi^i_j$,


определяемые этой леммой, называются формами связности корепера


$\tau$. Коэффициенты $\Gamma^i_{jk}$ в разложении


$\Phi^i_j= \Gamma^i_{jk}\tau^k$, $i,j,k=1,\dotsc,n$, называются


символами Кристоффеля. С их помощью по обычным формулам классического


тензорного анализа строится ковариантное дифференцирование. Например,


если $\xi=(\xi_i)_{i=1}^n$ --- локальный репер, дуальный кореперу


$\tau$, $A=A^i_j\xi_i\otimes\tau^j$ --- тензорное поле типа $(1,1)$,


то ковариантной производной его называется тензорное поле


$\nabla A=A^i_{j,k}\xi_i\otimes\tau^j\otimes\tau^k$


типа $(1,2)$, компоненты


которого $A^i_{j,k}$ определяются равенством


$A^i_{j,k}=A^i_{j|k}+\Gamma^i_{lk}A^l_j-\Gamma^l_{jk}A^i_l$,


где $A^i_{j|k}$ определяются


из разложения $dA^i_j=A^i_{j|k}\tau^k$, $i,j,k,l=1,\dotsc,n$. $\square$


\medskip





{\bf 1.5.} Через $\wedge^2T^*(X)$ будем обозначать расслоение


кососимметрических билинейных форм, через $S^2T^*(X)$ --- расслоение


симметрических билинейных форм. В расслоении $S^2T^*(X)$ определено


подрасслоение $S^2_{\Delta}T^*(X)$ невырожденных билинейных форм с


нормальным видом


\begin{equation}


ds^2=\sum_{i=1}^{n}\Delta^i\tau^i\otimes\tau^i.


\end{equation}


Сечения расслоения $S^2_{\Delta}T^*(X)$ называются метриками сигнатуры


$\Delta$. Поскольку метрика $ds^2$, определенная равенством


(1), не зависит от выбора локального корепера $\tau$ из


содержащей его орбиты, то всякая $\bold{SO(n,\bsy\Delta)}$-орбита


определяет некоторую метрику сигнатуры $\Delta$. Обратно, всякой


метрике $ds^2$ можно сопоставить две $\bold{SO(n,\bsy\Delta)}$-орбиты, в


кореперах из которых эта метрика имеет вид (1). Локальные


кореперы одной из этих орбит связаны с локальными кореперами другой


матрицами из $\bold{O(n,\bsy\Delta)}$ с определителями, равными $-1$. Будем


считать, что на $X$ зафиксирована ориентация,


и через $\Re_{\Delta}(X)$ обозначать


ту из этих двух орбит, локальные кореперы


которой связаны с координатным корепером $(dx^i)_{i=1}^n$ в карте из


ориентации матрицами с положительными определителями. О такой орбите


$\Re_{\Delta}(X)$ будем говорить, что она порождена метрикой


$ds^2$.


\medskip





{\bf 1.6.} Плоским $m$-мерным пространством, $m>n$, мы называем $m$-мерное


аффинное пространство $\Pi$, на векторном пространстве которого задана


невырожденная симметрическая билинейная форма, называемая скалярным


произведением векторов. Пусть $z:X\to\Pi$ --- $C^r$-погружение,


$r\ge 2$, $Tz$ --- касательное к $z$


отображение, $F=z(X)$ --- соответствующая поверхность в $\Pi$,


$TF$ --- касательное расслоение поверхности $F$, $T^{\bot}F$ ---


нормальное расслоение. Через $I(z)$ будем обозначать метрику на $X$,


индуцированную $z$. Если


$\Re_{\Delta}(X)$ --- $\bold{SO(n,\bsy\Delta)}$-орбита,


порожденная метрикой $I(z)$, то для всякого


локального корепера


$\tau=(\tau^i)_{i=1}^n\in{\cal E}^{r-1}(X,\Re_{\Delta}(X))$


определен локальный репер


$e=(e_i)_{i=1}^n\in{\cal E}^{r-1}(X,TF)$ по формулам $e_i=\xi_i(z)$, где


$\xi=(\xi_i)_{i=1}^n$ --- локальный репер на $X$, дуальный кореперу


$\tau$. Для векторов репера $e$ имеем


\begin{equation}


e_ie_j=\Delta_{ij}=


\begin{cases}


\Delta^i & \text{при }i=j,\\


0 & \text{при }i\ne j,


\end{cases}


\end{equation}


$i,j=1,\ldots,n$.





Будем предполагать, что в каждой точке $x\in X$ касательная плоскость


$T_xF$ поверхности $F$ содержит $n$ попарно ортогональных неизотропных


направлений. В этом случае в окрестности каждой точки $x\in X$


определен локальный $C^{r-1}$-репер $\nu=(\nu_{\sigma})_{\sigma=1}^p$,


$p=m-n$, в нормальном расслоении $T^{\bot}F$, удовлетворяющий условиям


\begin{equation}


\nu_{\sigma}\nu_{\tau}=


\begin{cases}


\Delta^{\sigma} & \text{при }\tau=\sigma,\\


0 & \text{при }\tau\ne\sigma,


\end{cases}


\end{equation}


$$


e_i\nu_{\sigma}=0,\quad i=1,\dotsc,n,\quad \sigma=1,\dotsc,p


$$


(\cite{Eis}, гл.\,4, \S\,42, c.\,175--179).


Мы пишем для краткости


$\Delta^{\sigma}$ вместо $\Delta^{n+\sigma}$.


Поскольку всякий вектор из $\Pi$ может быть представлен в виде


линейной комбинации векторов $e_i,\nu_{\sigma}$, $i=1,\dotsc,n$,


$\sigma=1,\dotsc,p$, сигнатура скалярного произведения в $\Pi$


имеет вид $\Delta=(e_1^2,\dotsc,e_n^2,\nu_1^2,\ldots,\nu_p^2)$. Об


аффинной системе координат $(O;a_1,\dotsc,a_m)$ в $\Pi$, базисные


векторы которой удовлетворяют соотношениям


$$


a_ia_j=e_ie_j,\quad a_ia_{n+\sigma}=0,\quad


a_{n+\sigma}a_{n+\tau}=\nu_{\sigma}\nu_{\tau},\quad


i,j=1,\dotsc,n,\quad \tau,\sigma=1,\dotsc,p,


$$


будем говорить, что она задает сигнатуру пространства $\Pi$.





Разложение дифференциалов $de_i$, $d\nu_{\sigma}$ по базису


$(e_i,\nu_{\sigma})$ приводит к формулам Гаусса и Вейнгартена


\begin{align}


de_i&=\Phi^k_ie_k+\sum_{\sigma=1}^{p}\Delta^{\sigma}\omega^{\sigma}_i


\nu_{\sigma},\notag\\


d\nu_{\sigma}&=-\sum_{i=1}^{n}\Delta^i\omega_i^{\sigma}e_i+


\sum_{\tau=1}^{p}\Delta^{\tau}\varkappa_{\sigma}^{\tau}\nu_{\tau},\\


i,k&=1,\dotsc,n,\quad \tau,\sigma=1,\dotsc,p,\quad p=m-n,\notag


\end{align}


где


$\Delta^k=e_k^2$, $\Delta^{\sigma}=\nu_{\sigma}^2$,


$\Phi^i_k$ --- формы связности локального корепера


$\tau\in {\cal E}^{r-1}(X,\Re_{\Delta}(X))$,


$\omega^{\sigma}_i$ --- 1-формы,


называемые формами погружения,


$\varkappa^{\sigma}_{\tau}= -\varkappa^{\tau}_{\sigma}$ --- 1-формы,


называемые формами кручения.


Формы погружения удовлетворяют соотношению


\begin{equation}


\omega^{\sigma}_i\wedge\tau^i=0,\quad i=1,\dotsc,n,\quad \sigma=1,\dotsc,p.


\end{equation}


Симметрическую билинейную форму


$$


II^{\sigma}=\omega^{\sigma}_i\otimes\tau^i=b^{\sigma}_{ij}


\tau^i\otimes\tau^j\in{\cal E}^{r-2}(X,S^2T^*(X))


$$


называют второй основной формой погружения $z$ (или поверхности $F$)


относительно нормали $\nu_{\sigma}$.





Если погружение $z\in C^3(X,\Pi)$, то внешнее дифференцирование формул


Гаусса и Вейнгартена приводит к уравнениям Гаусса, Петерсона-Кодацци и


Риччи


\begin{equation}


\begin{split}


&\sum_{\sigma=1}^{p}\Delta^{\sigma}\omega^{\sigma}_i\wedge\omega_j^


{\sigma}=\Delta^i(d\Phi^i_j+\Phi^i_k\wedge\Phi^k_j),\\


&d\omega_i^{\sigma}=\Phi^k_i\wedge\omega^{\sigma}_k+\sum_{\tau=1}^{p}


\Delta^{\tau}\omega^{\tau}_i\wedge\varkappa^{\sigma}_{\tau},\\


&d\varkappa^{\tau}_{\sigma}=\sum_{i=1}^{n}\Delta^i\omega_i^{\tau}\wedge


\omega_i^{\sigma}+\sum_{\rho=1}^{p}\Delta^{\rho}\varkappa_{\sigma}^{\rho}


\wedge\varkappa_{\rho}^{\tau},\\


&i,j,k=1,\dotsc,n,\quad \tau,\sigma=1,\dotsc,p.


\end{split}


\end{equation}





{\bf 1.7.} Цель данного и следующего пунктов --- получить уравнения Гаусса,


Петерсона-Кодацци и Риччи для $C^2$-погружений, используя обобщенное


внешнее дифференцирование по де~Раму. В этом пункте мы приведем


определение обобщенного внешнего дифференциала и некоторые его


свойства.





Пусть $U$ --- открытое множество на $X$ с компактным замыканием,


содержащимся в некоторой координатной окрестности. Обозначим через


$\wedge^lT^*(U)$ расслоение внешних дифференциальных форм степени


$l\ge 0$ на $U$, $l= 0,1,\dotsc,$ через


${\cal E}^{\infty}_0(U,\wedge^lT^*(U))$ ---


пространство $C^{\infty}$-сечений этого расслоения с


компактными носителями, содержащимися в $U$. Потоком степени $q$ на


$U$ называется всякий непрерывный линейный функционал


$v:{\cal E}^{\infty}_0(U,\wedge^{n-q}T^*(U))\to {\bold R}$


(\cite{DeRh}, гл.\,3, c.\,65--130). Границей $bv$ потока $v$


называется поток степени $q+1$, определяемый формулой


$$


bv[\varphi]=v[d\varphi] \ \forall \varphi\in


{\cal E}^{\infty}_0(U,\wedge^{n-q-1}T^*(U)).


$$


Каждая непрерывная $q$-форма $\omega\in {\cal E}^0(U,\wedge^qT^*(U))$


определяет поток $\Bar\omega$ степени $q$ на $U$ формулой


$$


\Bar\omega[\varphi]=\int_{U}\omega\wedge\varphi \


\forall \varphi\in{\cal E}^{\infty}_0(U,\wedge^{n-q}T^*(U)).


$$





Непрерывную форму $\Omega\in {\cal E}^0(X,\wedge^{q+1}T^*(X))$


будем называть обобщенным внешним дифференциалом непрерывной формы


$\omega\in {\cal E}^0(X,\wedge^qT^*(X))$ на $U$, если


$\overline\Omega= (-1)^{q+1}b\Bar\omega$, т.\,е., если


\begin{equation}


\int_{U}\Omega\wedge\varphi=(-1)^{q+1}\int_{U}\omega\wedge d\varphi \


\forall \varphi\in{\cal E}^{\infty}_0(U,\wedge^{n-q-1}T^*(U))


\end{equation}


(ср. \cite{Bor1}--\cite{Bor3}). Обозначим через


${\cal E}^1_0(U,\wedge^lT^*(U))$ множество всех $C^1$-сечений


расслоения $\wedge^lT^*(U)$ с компактными носителями, содержащимися


в $U$. Так как пространство ${\cal E}^{\infty}_0(U,\wedge^lT^*(U))$


всюду плотно в ${\cal E}^1_0(U,\wedge^lT^*(U))$, то $\Omega$


является обобщенным внешним дифференциалом формы $\omega$ тогда и


только тогда, когда равенство (7) имеет место для всякой формы


$\varphi\in {\cal E}^1_0(U,\wedge^{n-q-1}T^*(U))$. Отсюда следует,


что если $\omega\in {\cal E}^1(U,\wedge^qT^*(U))$, то обобщенный


внешний дифференциал $d\omega$ совпадает с обычным внешним


дифференциалом. Кроме того, $d\omega$ (если он существует)


определяется формой $\omega$ однозначно, не зависит от системы


координат на $U$, имеет локальный характер, и $dd\omega= 0$.





Следующие два предложения касаются свойств обобщенного внешнего


дифференциала. Доказательства их приведены в \cite{Mark4}.





\begin{lem}


Если форма $\omega\in {\cal E}^0(X,\wedge^qT^*(X))$ имеет


обобщенный внешний дифференциал $d\omega$, то для всякой $p$-формы


$\psi\in {\cal E}^1(X,\wedge^pT^*(X))$, $p,q\ge0$, справедливо


равенство


$d(\omega\wedge \psi)= d\omega\wedge \psi+ (-1)^q\omega\wedge d\psi$.


\end{lem}





\begin{lem}


Если многообразие $X$ односвязно, то для всякой формы


$\omega\in {\cal E}^0(X,T^*(X))$, для которой $d\omega= 0$, найдется


функция $f\in C^1(X,{\bold R})$, такая, что $\omega= df$.


\end{lem}





{\bf 1.8.} В дальнейшем нам понадобится


\begin{thm}


Для всякого $C^2$-погружения $z: X\to \Pi$ формы связности


$\Phi^i_j$, погружения $\omega^{\sigma}_i$ и кручения


$\varkappa^{\sigma}_{\tau}$, $i,j= 1,\dotsc,n$,


$\tau,\sigma= 1,\dotsc,p$, обладают обобщенным внешним дифференциалом


и удовлетворяют уравнениям Гаусса, Петерсона-Кодацци и Риччи


{\em (6)}, в которых $d$ --- знак обобщенного внешнего дифференциала.


\end{thm}





\begin{pf} Пусть $U$ --- окрестность произвольной точки


$x\in X$ с компактным замыканием, содержащимся в некоторой


координатной окрестности. Если $z\in C^2(X,\Pi)$, то


$de_i\in {\cal E}^0(U,T^*(U)\otimes \Pi)$,


$i= 1,\dotsc,n$. В силу леммы 3,


учитывая, что $dde_i= 0$, для всякой формы


$\varphi\in {\cal E}^{\infty}_0(U,\wedge^{n-2}T^*(U))$ имеем


$d(de_i\wedge \varphi)= -de_i\wedge d\varphi$. Отсюда следует, что


$\int_U de_i\wedge d\varphi= 0$. В силу формул Гаусса (4), это


равенство влечет


$$


\int_{U}\Bigl(\Phi^k_ie_k+\sum_{\sigma=1}^{p}\Delta^{\sigma}\omega^{\sigma}_i


\nu_{\sigma}\Bigr)\wedge d\varphi=0,\quad i,k=1,\dotsc,n.


$$


Поскольку ${\cal E}^{\infty}_0(U,\wedge^{n-2}T^*(U))$ всюду плотно в


${\cal E}^1_0(U,\wedge^{n-2}T^*(U))$, последнее равенство


справедливо и для всякой формы


$\varphi\in {\cal E}^1_0(U,\wedge^{n-2}T^*(U))$. Внося $e_k$ и


$\nu_{\sigma}$ под знак $d$, пользуясь леммой 3, получаем


$$


\int_{U}\biggl[\Phi^k_i\wedge d(\varphi e_k)+


\sum_{\sigma=1}^{p}\Delta^{\sigma}\omega^{\sigma}_i\wedge


d(\varphi\nu_{\sigma})-\Phi^k_i\wedge de_k\wedge\varphi-


\sum_{\sigma=1}^{p}\Delta^{\sigma}\omega^{\sigma}_i\wedge


d\nu_{\sigma}\wedge\varphi\biggr]=0.


$$


В силу (4) отсюда находим


\begin{gather*}


\int_{U}\biggl[\Phi^k_i\wedge d(\varphi e_k)+


\sum_{\tau=1}^{p}\Delta^{\tau}\omega^{\tau}_i\wedge


d(\varphi\nu_{\tau})-


\biggl(\Phi^k_i\wedge\Phi^j_k


-\sum_{\sigma=1}^{p}


\Delta^{\sigma}\omega^{\sigma}_i\wedge\omega^\sigma_j


\biggr)\wedge(e_j\varphi)-\\


-\sum_{\tau=1}^{p}\Delta^{\tau}\biggl(\Phi^k_i\wedge\omega^{\tau}_k+


\sum_{\sigma=1}^{p}\Delta^{\sigma}\omega^{\sigma}_i\wedge


\varkappa^{\tau}_{\sigma}\biggr)\wedge(\nu_{\tau}\varphi)\biggr]=0.


\end{gather*}





Пусть $e_i= e^{\alpha}_ia_{\alpha}$,


$\nu_{\sigma}=\nu_{\sigma}^{\alpha}a_{\alpha}$,


$\alpha= 1,\dotsc,m$, $i= 1,\dotsc,n$, $\sigma= 1,\dotsc,p$, где


$(O,a_{\alpha})_{\alpha=1}^m$ --- аффинная система координат в $\Pi$,


задающая сигнатуру. Полагая в последнем интегральном равенстве


$\varphi= e^{\alpha}_i\psi$,


$\psi\in {\cal E}^1_0(U,\wedge^{n-2}T^*(U))$,


умножая его скалярно на $a_{\alpha}$,


после суммирования по $\alpha= 1,\ldots,m$ получим


\begin{equation}


\int_{U}\Phi^j_i\wedge d\psi=\int_{U}\biggl(\Phi^k_i\wedge\Phi^j_k-


\sum_{\sigma=1}^{p}\Delta^i\Delta^{\sigma}\omega^{\sigma}_i\wedge


\omega^{\sigma}_j\biggr)\wedge\psi.


\end{equation}


Если же положить $\varphi= \nu^{\alpha}_{\sigma}\psi$, то таким же


образом получим


\begin{equation}


\int_{U}\Delta^{\tau}\omega^{\tau}_i\wedge d\psi=


\int_{U}\biggl(\Delta^{\tau}\Phi^k_i\wedge\omega^{\tau}_k+


\sum_{\sigma=1}^{p}\Delta^{\sigma}\omega^{\sigma}_i\wedge


\varkappa^{\tau}_{\sigma}\biggr)\wedge\psi.


\end{equation}


Поскольку равенства (8) и (9) справедливы для всякой


формы


$\psi\in {\cal E}^1_0(U,\wedge^{n-2}T^*(U))$, в силу определения


обобщенного внешнего дифференциала, получаем, что формы $\Phi^i_j$ и


$\omega^{\sigma}_i$, $i,j= 1,\dotsc,n$, $\sigma= 1,\dotsc,p$, обладают


обобщенными внешними дифференциалами, совпадающими с правыми частями


соответствующих уравнений (6) в каждой точке $x\in X$.





Рассмотрим тождество


$d(d\nu_{\sigma}\wedge \varphi)= -d\nu_{\sigma}\wedge d\varphi$,


$\sigma= 1,\dotsc,p$. Из него следует, что


$\int\limits_{U}d\nu_{\sigma}\wedge d\varphi= 0$. В силу (4), получаем


$$


\int_{U}\biggl[-\sum_{i=1}^{n}\Delta^i\omega^{\sigma}_ie_i+


\sum_{\tau=1}^{p}\Delta^{\tau}\varkappa^{\tau}_{\sigma}\nu_{\tau}\biggr]


\wedge d\varphi=0,\quad \tau,\sigma=1,\dotsc,p,\quad i=1,\dotsc,n.


$$


Пользуясь леммой 3, находим


\begin{gather*}


\int_{U}\biggl[-\sum_{i=1}^{n}\Delta^i\omega^{\sigma}_i


\wedge d(\varphi e_i)+


\sum_{\tau=1}^{p}\Delta^{\tau}\varkappa^{\tau}_{\sigma}\wedge


d(\varphi\nu_{\tau})+\\


+\biggl(\sum_{i=1}^{n}\Delta^i\omega^{\sigma}_i\wedge de_i-


\sum_{\tau=1}^{p}\Delta^{\tau}\varkappa^{\tau}_{\sigma}\wedge


d\nu_{\tau}\biggr)\wedge\varphi\biggr]=0.


\end{gather*}


В силу (4) отсюда имеем


\begin{gather*}


\int_{U}\biggl[-\sum_{i=1}^{n}\Delta^i\omega^{\sigma}_i\wedge d(\varphi e_i)+


\sum_{\tau=1}^{p}\Delta^{\tau}\varkappa^{\tau}_{\sigma}\wedge


d(\varphi\nu_{\tau})+\\


+\sum_{i=1}^{n}\Delta^i\omega^{\sigma}_i\wedge\biggl(\Phi^k_ie_k+


\sum_{\rho=1}^{p}\Delta^{\rho}\omega^{\rho}_i\nu_{\rho}\biggr)


\wedge\varphi-\\


-\sum_{\tau=1}^{p}\Delta^{\tau}\varkappa^{\tau}_{\sigma}


\wedge\biggl(-\sum_{i=1}^{n}\Delta^i\omega^{\tau}_ie_i+


\sum_{\rho=1}^{p}\Delta^{\rho}\varkappa^{\rho}_{\tau}\nu_{\rho}


\biggr)\wedge\varphi\biggr]=0.


\end{gather*}


Полагая $\varphi= \nu^{\alpha}_{\rho}\psi$, умножая на $a_{\alpha}$,


после суммирования по $\alpha= 1,\ldots,m$ получаем


$$


\int_{U}\Delta^{\rho}\varkappa^{\rho}_{\sigma}\wedge d\psi=


-\int_{U}\biggl(\,\sum_{i=1}^{n}\Delta^{\rho}\Delta^i


\omega^{\sigma}_i\wedge\omega^{\rho}_i-


\sum_{\tau=1}^{p}\Delta^{\rho}\Delta^{\tau}\varkappa^{\tau}_{\sigma}


\wedge\varkappa^{\rho}_{\tau}\biggr)\wedge\psi.


$$


Отсюда и из определения обобщенного внешнего дифференциала, в силу


произвольности $\psi\in {\cal E}^1_0(U,\wedge^{n-2}T^*(U))$


следует, что формы $\varkappa^{\rho}_{\sigma}$,


$\rho,\sigma= 1,\ldots,p$, обладают обобщенным внешним дифференциалом


в точке $x\in X$, и этот дифференциал совпадает с правой частью


соответствующего равенства из (6).


\end{pf}





\section{Поля вращений}





{\bf 2.1.} Пусть $V(X)$ --- векторное расслоение, $f\in {\cal E}^r(X,V(X))$,


$r\ge0$. Непрерывной деформацией сечения $f$ будем называть всякое


непрерывное отображение


$F: X\times (-\varepsilon,\varepsilon)\to V(X)$, $\varepsilon> 0$,


удовлетворяющее условиям


\begin{itemize}


\item[{\bf a)}] $F(x,0)=f(x)$ {\it для всякой точки} $x\in X$;\\


\item[{\bf б)}] {\it для всякого $t\in (-\varepsilon,\varepsilon)$ отображение


$f_t: X\to V(X)$, определенное равенством $f_t(x)= F(x,t)$,


является сечением из} ${\cal E}^r(X,V(X))$.


\end{itemize}





Будем говорить, что деформация $F$ принадлежит классу $C^s$, $s\ge1$,


если


\begin{itemize}


\item[{\bf в)}] {\it для всякой точки $x\in X$ отображение


$F_x: (-\varepsilon,\varepsilon)\to V(X)$, определенное равенством


$F_x(t)=F(x,t)$, принадлежит классу} $C^s$;\\


\item[{\bf г)}] {\it дифференцирование функции $F(x,t)$ по


$t\in (-\varepsilon,\varepsilon)$ до порядка $s$ включительно


перестановочно с дифференцированием по локальным координатам точки $x$


до порядка $r$ включительно.}


\end{itemize}





Говорят, что деформация $F$ принадлежит классу $C^{\infty}$, если она


принадлежит классу $C^s$ для всех $s= 1,2,\dotsc$ Деформация $F$


называется аналитической, если она принадлежит классу $C^{\infty}$ и в


каждой точке $x\in X$ может быть представлена степенным рядом


\begin{equation}


F(x,t)=f(x)+2\sum_{s=1}^{\infty}\delta^sf(x)t^s,


\end{equation}


сходящимся на $(-\varepsilon,\varepsilon)$, коэффициенты которого


$\delta^sf\in {\cal E}^r(X,V^{(s)}(X))$.





Пусть $P(X)$ --- $C^r$-подрасслоение векторного расслоения $V(X)$,


$f\in{\cal E}^r(X,P(X))$. Деформацией класса $C^s$,


$1\le s\le \infty$, сечения $f$ будем называть всякое непрерывное


отображение


$F: X\times (-\varepsilon,\varepsilon)\to P(X)$, $\varepsilon> 0$,


удовлетворяющее условиям а) и б) и являющееся деформацией $f$ класса


$C^s$ как сечения из ${\cal E}^r(X,V(X))$. Деформацию сечения


$f\in{\cal E}^r(X,P(X))$ будем называть аналитической, если она


является деформацией класса $C^{\infty}$ и в каждой точке может быть


представлена сходящимся на $(-\varepsilon,\varepsilon)$ рядом


(10), в котором $\delta^sf\in {\cal E}^r(X,P^{(s)}(X))$.


Каноническое вложение $P_x(X)\subset V_x(X)$ позволяет в каждой точке


$x\in X$ значение $\delta^{(s)}f(x)$, $s=1,2,\dotsc$,


рассматривать как вектор из $V_x(X)$.





Использование $s$-производного расслоения в определении аналитической


деформации обусловлено тем, что для фиксированной точки $x\in X$


аналитическая деформация $F$ задает аналитический путь


$F_x: (-\varepsilon,\varepsilon)\to P_x(X)\subset V_x(X)$. При этом


\begin{equation}


\left.\delta^sf(x)=\frac{d^sF_x(t)}{2s!dt^s}\right|_{t=0}, 


s=1,2,\dotsc


\end{equation}


Для заданной деформации $F$ класса $C^s$ сечение


$\delta^sf\in {\cal E}^r(X,P^{(s)}(X))$, определяемое формулой


(11), называется


$s$-й вариацией сечения $f$. Полагаем $\delta^0f= f$,


$\delta^1f= \delta f$.





Класс деформаций, определяющих одни и те же вариации


$\delta^1f,\dotsc,\delta^lf$, называется б.\,м. деформацией $l$-го


порядка сечения $f$ (ср. \cite{Mark5}). Б.\,м. деформация $l$-го


порядка однозначно определяется заданием какой-либо деформации класса


$C^l$ с заданными вариациями $\delta^1f,\dotsc,\delta^lf$. Всякая


такая деформация может быть записана в виде


$F(x,t)=f(x)+2\sum\limits_{s=1}^{l}\delta^sf(x)t^s+o(l)$,


где $o(l)\in{\cal E}^r(X,V(X))$, $\dfrac{o(l)}{t}\to 0$ при $t\to 0$. Тем


самым б.\,м. деформацию $l$-го порядка можно рассматривать как $l$-струю


\cite{Kli} аналитической деформации. Оператор


$\delta^s: {\cal E}^r(X,P(X))\to {\cal E}^r(X,P^{(s)}(X))$,


определяемый формулой (11), называется оператором варьирования


$s$-го порядка.


\medskip





{\bf 2.2.} Рассмотрим б.\,м. деформацию


$\{z^t\}_{t\in(-\varepsilon,\varepsilon)}$


$l$-го порядка погружения $z: X\to \Pi$, $l= 1,\dotsc,\infty$.


При $l= \infty$ деформацию считаем аналитической. Рассматривая $z$ как


сечение из ${\cal E}^r(X,X\times \Pi)$, $r\ge 1$, мы можем


пользоваться терминологией предыдущего пункта. Деформация


$\{z^t\}_{t\in(-\varepsilon,\varepsilon)}$ порождает деформацию


$\{I(z^t)\}_{t\in(-\varepsilon,\varepsilon)}$ метрики $I(z)$, а значит,


и деформацию


$\tau_t= (\tau^i_t)_{i=1}^n$, $t\in (-\varepsilon,\varepsilon)$,


локального корепера $\tau\in {\cal E}^{r-1}(X,\Re_{\Delta}(X))$,


где $\tau_t\in {\cal E}^{r-1}(X,\Re_{\Delta}^t(X))$,


$\Re_{\Delta}^t(X)$ ---


$\bold{SO(n,\bsy\Delta)}$-орбита метрики $I(z^t)$. При


этом деформация $\{\tau_t\}_{t\in (-\varepsilon,\varepsilon)}$


определяется с точностью до наложения деформации корепера $\tau$ по


орбите $\Re_{\Delta}(X)$, т.\,е. деформации вида


$\tau_t^i= q^i_j(t)\tau^j$, где матрица


$(q^i_j(t))_{i,j=1}^n\in C^{r-1}(X,{\bold{SO(n,\bsy\Delta)}})$. Отсюда


следует, что деформация локального репера $e= (e_j)_{j=1}^n$ из


касательного расслоения $TF$ определяется деформацией погружения $z$ c


точностью до наложения деформации вида $e^t_j= p^i_j(t)e_i$, где


$(p^i_j(t))_{i,j=1}^n\in C^{r-1}(X,{\bold{SO(n,\bsy\Delta)}})$. Из


(3) следует, что деформация локального репера


$\nu= (\nu_{\sigma})_{\sigma=1}^p$ нормального расслоения $T^{\bot}F$


определяется с точностью до наложения деформации вида


$\nu^t_{\sigma}= q^{\tau}_{\sigma}(t)\nu_{\tau}$, где


$(q^{\tau}_{\sigma}(t))_{\tau,\sigma=1}^p\in C^{r-1}(X,{\bold{


SO(p,\bsy\Delta)}})$.





Будем считать, что при заданной деформации


$\{z^t\}_{t\in(-\varepsilon,\varepsilon)}$ деформации реперов $e$ и


$\nu$ зафиксированы, и называть их порожденными деформацией


$\{z^t\}_{t\in(-\varepsilon,\varepsilon)}$.





\begin{thm}


Для заданной б.\,м. деформации $l$-го порядка $z^t: X\to \Pi$,


$t\in (-\varepsilon,\varepsilon)$, $\varepsilon> 0$,


$l= 1,\dotsc,\infty$, погружения $z\in C^r(X,\Pi)$, $r\ge 1$, и


порожденных ею деформаций локальных реперов $e= (e_i)_{i=1}^n$,


$\nu= (\nu_{\sigma})_{\sigma=1}^p$ существует и притом единственная


система бивекторов $(V^s)_{s=1}^l$, $V^s\in C^{r-1}(X,\wedge^2\Pi)$,


такая, что


\begin{equation}


\delta^se_i=\sum_{\alpha=1}^{s}V^{\alpha}\delta^{s-\alpha}e_i,\quad


\delta^s\nu_{\sigma}=\sum_{\alpha=1}^{s}V^{\alpha}\delta^{s-\alpha}


\nu_{\sigma},\quad


s=1,\dotsc,l, \ i=1,\dotsc,n, \sigma=1,\dotsc,p,


\end{equation}


где под знаками суммы --- внутреннее произведение бивектора на вектор


{\em(\cite{Car2}, ч.\,1, \S\,2, c.\,15).}


\end{thm}





В доказательстве этой теоремы мы используем следующее элементарное


предложение линейной алгебры, доказательство которого опускаем.





\begin{lem}


Если для систем векторов $b_{\alpha}= b^{\lambda}_{\alpha}a_{\lambda}$


и $c_{\beta}= c^{\mu}_{\beta}a_{\mu}$,


$\alpha,\beta,\lambda,\mu= 1,\dotsc,m$ в $\Pi$ выполнены условия


$c_{\alpha}c_{\beta}= b_{\alpha}b_{\beta}$,


$\det(b^{\lambda}_{\alpha})\det(c^{\mu}_{\beta})> 0$, то существует и


притом единственная матрица


$P= (p^{\alpha}_{\beta})_{\alpha,\beta=1}^m\in{\bold{SO(m,\bsy\Delta)}}$,


такая, что $c_{\beta}= p^{\alpha}_{\beta}b_{\alpha}$.


\end{lem}





Докажем теорему 2. Для векторов реперов $e^t= (e^t_i)_{i=1}^n$


и $\nu^t= (\nu^t_{\sigma})_{\sigma=1}^p$, порожденных деформацией


$z^t: X\to \Pi$, в каждой точке $x\in X$ и при каждом


$t\in (-\varepsilon,\varepsilon)$ имеем


$$


e^t_ie^t_j=\Delta_{ij},\ e^t_i\nu^t_{\sigma}=0, \ \nu^t_{\sigma}


\nu^t_{\tau}=\Delta_{\sigma\tau}, \ i,j=1,\dotsc,n, \ \sigma,


\tau=1,\dotsc,p.


$$


Отсюда и из (2) и (3), в силу леммы 5,


следует существование матрицы


$P(t)= (p^{\alpha}_{\beta}(t))_{\alpha,\beta=1}^m\in {\bold{


SO(m,\bsy\Delta)}}$, для которой


\begin{equation}


e^t_i=P(t)e_i,\ \nu^t_{\sigma}=P(t)\nu_{\sigma},\ i=1,\dotsc,n, 


\sigma=1,\dotsc,p,\ t\in(-\varepsilon,\varepsilon),


\end{equation}


где в правых частях --- произведения матриц, и отображение


$P: (-\varepsilon,\varepsilon)\to {\bold{SO(m,\bsy\Delta)}}$ является


аналитическим путем в $\bold{SO(m,\bsy\Delta)}$ с $P(0)= 1$ в каждой точке


$x\in X$. По лемме 1 можем положить $P(t)= \theta(v(t))$, где


$v: (-\varepsilon,\varepsilon)\to \wedge(m)$ --- аналитический


путь в пространстве $\wedge(m)$ кососимметрических $m\times m$-матриц


с $v(0)= 0$. Отсюда для вариации $\delta^sP$ в силу (11)


находим


\begin{equation}


\begin{split}


\delta^sP&=\frac{v^{(s)}(0)}{s!}\Delta+\sum_{\alpha=1}^{s-1}


\frac{v^{(\alpha)}(0)}{\alpha!}\delta^{s-\alpha}P,\\


\delta P&=v^{(1)}(0)\Delta,\quad s=2,\dotsc,l.


\end{split}


\end{equation}





Пусть $v^{(s)}(0)= (v^{\lambda\mu}_s)_{\lambda,\mu=1}^m$,


$e_i= e_i^{\lambda}a_{\lambda}$,


$\nu_{\sigma}= \nu_{\sigma}^{\lambda}a_{\lambda}$,


$\lambda= 1,\dotsc,m$, $i= 1,\dotsc,n$, $\sigma=1,\dotsc,p$. Из


(13) и (14) находим


\begin{align}


\delta^se_i^{\mu}&=\frac{v_{s}^{\lambda\mu}}{s!}\Delta_{\lambda\rho}


e_i^{\rho}+


\sum_{\alpha=1}^{s-1}\frac{v_{\alpha}^{\lambda\mu}}{\alpha!}


\Delta_{\lambda\rho}\delta^{s-\alpha}e_i^{\rho},\notag\\


\delta^s\nu_{\sigma}^{\mu}&=\frac{v_{s}^{\lambda\mu}}{s!}\Delta_{\lambda\rho}


\nu_{\sigma}^{\rho}+


\sum_{\alpha=1}^{s-1}\frac{v_{\alpha}^{\lambda\mu}}{\alpha!}


\Delta_{\lambda\rho}\delta^{s-\alpha}\nu_{\sigma}^{\rho},\quad


s=2,\dotsc,l,\\


\delta e^{\mu}_i&=v^{\lambda\mu}_1\Delta_{\lambda\rho}e^{\rho}_i,\quad


\delta\nu_{\sigma}^{\mu}=v^{\lambda\mu}_1\Delta_{\lambda\mu}


\nu^{\rho}_{\sigma},\notag


\end{align}


$$


i=1,\dotsc,n,\ \sigma=1,\dotsc,p,\ \lambda,\mu,\rho,=1,\dotsc,m .


$$


Рассмотрим в $\wedge^2\Pi$ систему бивекторов


$$


V^s=\frac{1}{2s!}v^{\lambda\mu}_s[a_{\lambda},a_{\mu}], \


s=1,\dotsc,l,\ \lambda,\mu=1,\dotsc,m,


$$


где $[*,*]$ --- внешнее умножение векторов в $\Pi$. С использованием


бивекторов $V^s$ равенства (15) запишутся в виде (12).


Тем самым доказано существование системы $(V^s)_{s=1}^l$.





Докажем единственность. Допустим, что существует еще одна система


полей $(\widetilde V^s)_{s=1}^l$, удовлетворяющая условиям (12).


Для разности $\widetilde V^1-V^1$ имеем


$$


(\widetilde V^1-V^1)e_i=0,\quad (\tilde V^1-V^1)\nu_{\sigma}=0,\quad


i=1,\dotsc,n, \quad


\sigma=1,\dotsc,p.


$$


Так как $(e_i,\nu_{\sigma})$ --- базис в $\Pi$, то получаем


$\widetilde V^1= V^1$. Если теперь по индукции допустить, что


$\widetilde V^1= V^1,\dotsc,\widetilde V^{s-1}= V^{s-1}$, то для разности


$\widetilde V^s- V^s$ из (12) снова получим


$(\widetilde V^s- V^s)e_i= 0$, $(\widetilde V^s- V^s)\nu_{\sigma}= 0$,


и, значит, $\widetilde V^s= V^s$ для всех $s= 1,\dotsc,l$.





Равенства (12) в каждой точке $x\in X$ и при каждом


$s= 1,\dotsc,l$ являются линейными относительно компонентов поля


$V^s$. Следовательно, эти компоненты являются рациональными функциями


от компонентов векторов $e_i,\nu_{\sigma}$ и бивекторов


$V^1,\dotsc,V^{s-1}$, причем компоненты $V^1$ зависят только от


компонентов $e_i$ и $\nu_{\sigma}$. Так как поля


$e_i,\nu_{\sigma}\in C^{r-1}(X,\Pi)$, то отсюда следует, что


$V^s\in C^{r-1}(X,\wedge^2\Pi)$, $s=1,\dotsc,l$. $\square$





Систему полей $(V^s)_{s=1}^l$, определяемую теоремой 2, будем


называть системой полей вращений, а ее элемент $V^s$ --- полем


вращений порядка $s$ при б.\,м. деформации $l$-го порядка (аналитической


деформации, если $l= \infty$) погружения $z: X\to \Pi$.


\medskip





{\bf 2.3.} Тот факт, что введенная система полей вращений для общей


б.м. деформации $l$-го порядка обобщает известную систему полей


вращений для б.м. изгибания $l$-го порядка, иллюстрируется следующей


теоремой.





\begin{thm}


Для системы полей вращений $(V^s)_{s=1}^l$ б.\,м. деформации


$z^t: X\to \Pi$, $t\in (-\varepsilon,\varepsilon)$,


$\varepsilon> 0$, $C^r$-погружения $z: X\to \Pi$, $r\ge 2$,


справедливы равенства


\begin{equation}


d\delta^sz=\sum_{\alpha=1}^{s}V^{\alpha}d\delta^{s-\alpha}z+


\sum_{\alpha=1}^{s-1}V^{\alpha}\delta^{s-\alpha}\tau^ie_i+


\delta^s\tau^ie_i,\quad


i=1,\dotsc,n,\ s=1,\dotsc,l,\ l=1,\dotsc,\infty


\end{equation}


{\em (}при $s=1$ вторая сумма исчезает{\em )}.


\end{thm}





\begin{pf} Имеем


\begin{align*}


d\delta^sz &=\delta^s(\tau^ie_i)=


\sum_{\alpha=1}^{s}(\delta^{\alpha}\tau^i\delta^{s-\alpha}e_i+


\delta^{s-\alpha}\tau^i\delta^{\alpha}e_i)=\\


&=\delta^s\tau^ie_i+\tau^i\delta^se_i+


2\sum_{\alpha=1}^{s-1}\delta^{\alpha}\tau^i\delta^{s-\alpha}e_i.


\end{align*}


Последнюю сумму обозначим через $M$. По теореме 2 она


преобразуется следующим образом:


\begin{align*}


M&=2\sum_{\alpha=1}^{s-1}\delta^{\alpha}\tau^i\sum_{\beta=1}^{s-\alpha}


V^{\beta}\delta^{s-\alpha-\beta}e_i=


2\sum_{\alpha=1}^{s-1}V^{\alpha}\sum_{\beta=1}^{s-\alpha}\delta^{\beta}


\tau^i\delta^{s-\alpha-\beta}e_i=\\


&=\sum_{\alpha=1}^{s-1}V^{\alpha}\biggl[\,\sum_{\beta=1}^{s-\alpha}(


\delta^{\beta}\tau^i\delta^{s-\alpha-\beta}e_i+


\delta^{s-\alpha-\beta}\tau^i\delta^{\beta}e_i)+


\delta^{s-\alpha}\tau^ie_i-\tau^i\delta^{s-\alpha}e_i\biggr]=\\


&=\sum_{\alpha=1}^{s-1}V^{\alpha}d\delta^{s-\alpha}z-\tau^i


\sum_{\alpha=1}^{s-1}V^{\alpha}\delta^{s-\alpha}e_i+


\sum_{\alpha=1}^{s-1}V^{\alpha}\delta^{s-\alpha}\tau^ie_i=\\


&=\sum_{\alpha=1}^{s}V^{\alpha}d\delta^{s-\alpha}z+\tau^i\delta^se_i+


\sum_{\alpha=1}^{s-1}V^{\alpha}\delta^{s-\alpha}\tau^ie_i.


\end{align*}


Отсюда вытекает требуемое равенство.


\end{pf}





В случае б.\,м. изгибания $l$-го порядка можем считать, что


$\delta^s\tau^i= 0$, $i= 1,\dotsc,n$, $s= 1,\dotsc,l$, и равенство


(16) превращается в известное равенство, определяющее поля


вращений б.\,м. изгибания \cite{Mark2}.


\medskip





{\bf 2.4.} В этом пункте мы покажем, что дифференциал поля вращений $s$-го


порядка выражается через $s$-е вариации форм связности, погружения и


кручения, а также через поля вращений и вариации указанных форм


порядков, меньших $s$.





\begin{lem}


Для дифференциалов полей вращений при б.\,м. деформации $l$-го порядка,


$l= 1,\dotsc,\infty$, погружения $z\in C^2(X,\Pi)$ имеют место


разложения


\begin{equation}


dV=-\frac{1}{2}\Delta^j\delta\Phi^i_j[e_i,e_j]+\Delta^i\Delta^{\sigma}


\delta\omega^{\sigma}_i[e_i,\nu_{\sigma}]-\frac{1}{2}\Delta^{\tau}


\Delta^{\sigma}\delta\varkappa^{\tau}_{\sigma}


[\nu_{\tau},\nu_{\sigma}],


\end{equation}


\begin{equation}


dV^s =\frac{1}{2}A^{ij}_s[e_i,e_j]+B^{i\sigma}_s[e_i,\nu_{\sigma}]


+\frac{1}{2}C^{\tau\sigma}_s[\nu_{\tau},\nu_{\sigma}], s\ge2,


\end{equation}


где $V\equiv V^1$, по индексам $i,j= 1,\dotsc,n$,


$\tau,\sigma= 1,\dotsc,p$ проводится суммирование, и


\begin{equation}


\begin{split}


A^{ij}_s &=-A^{ji}_s=\\


&=\frac{1}{2}\Delta^i\Delta^j\sum_{\alpha=1}^{s-1}\{


V^{\alpha}([\delta^{s-\alpha}\Phi^k_ie_k+


\Delta^{\sigma}\delta^{s-\alpha}\omega^{\sigma}_i\nu_{\sigma},e_j]+\\


&+[e_i,\delta^{s-\alpha}\Phi^k_je_k+


\Delta^\sigma\delta^{s-\alpha}\omega^{\sigma}_j\nu_{\sigma}])-\\


&-dV^{\alpha}([\delta^{s-\alpha}e_i,e_j]+


[e_i,\delta^{s-\alpha}e_j])\}+\Delta^i\delta^s\Phi^j_i,\\[2mm]


%


B^{i\sigma}_s &= \Delta^i\Delta^{\sigma}\sum_{\alpha=1}^{s-1}\{


V^{\alpha}[\delta^{s-\alpha}\Phi^k_ie_k+


\Delta^{\tau}\delta^{s-\alpha}\omega^{\tau}_i\nu_{\tau},\nu_{\sigma}]-\\


& - dV^{\alpha}[\delta^{s-\alpha}e_i,\nu_{\sigma}]\}+


\Delta^i\Delta^{\sigma}\delta^s\omega^{\sigma}_i,\\[2mm]


%


C^{\tau\sigma}_s &= -C^{\sigma\tau}_s=\\


& =\frac{1}{2}\Delta^{\sigma}\Delta^{\tau}\sum_{\alpha=1}^{s-1}\{


V^{\alpha}([-\Delta^i\delta^{s-\alpha}\omega^{\tau}_ie_i+


\Delta^{\rho}\delta^{s-\alpha}\varkappa^{\rho}_{\tau}


\nu_{\rho},\nu_{\sigma}]+\\


&+[\nu_{\tau},-\Delta^i\delta^{s-\alpha}\omega^{\sigma}_ie_i+


\Delta^{\rho}\delta^{s-\alpha}\varkappa^{\rho}_{\sigma}\nu_{\rho}])-\\


&-dV^{\alpha}([\delta^{s-\alpha}\nu_{\tau},\nu_{\sigma}]+


[\nu_{\tau},\delta^{s-\alpha}\nu_{\sigma}])\}


+\Delta^{\sigma}\Delta^{\tau}\delta^s\varkappa_{\tau}^{\sigma},\\


& k=1,\dotsc,n,\ \rho=1,\dotsc,p,\ s=1,\dotsc,l.


\end{split}


\end{equation}


\end{lem}





\begin{pf} Поскольку в каждой точке $x\in X$ бивекторы


$[e_i,e_j]$,


$[e_i,\nu_{\sigma}]$, $[\nu_{\tau},\nu_{\sigma}]$,


при $i<j$, $\tau<\sigma$, $i,j=1,\dotsc,n$,


$\tau,\sigma=1,\dotsc,p$, образуют базис пространства


$\wedge^2\Pi$, разложение (18) имеет место, и


можем считать, что $A^{ij}_s= -A^{ji}_s$, $C_s^{\tau\sigma}= 


-С^{\sigma\tau}_s$. Остается показать, что коэффициенты этого


разложения имеют вид


(19). Из (18) находим


\begin{gather}


A^{ij}_s=\Delta^i\Delta^jdV^s[e_i,e_j], \quad


B^{i\sigma}_s=\Delta^i\Delta^{\sigma}dV^s[e_i,\nu_{\sigma}],\notag\\


C^{\tau\sigma}_s=\Delta^{\tau}\Delta^{\sigma}dV^s


[\nu_{\tau},\nu_{\sigma}],\\


i,j=1,\dotsc,n,\ \tau,\sigma=1,\dotsc,p,\ s=1,\dotsc,l.\notag


\end{gather}


Дифференцируя первое из равенств (12), будем иметь


\begin{equation}


d\delta^se_i=\sum_{\alpha=1}^{s}dV^{\alpha}\delta^{s-\alpha}e_i+


\sum_{\alpha=1}^{s}V^{\alpha}d\delta^{s-\alpha}e_i.


\end{equation}


Преобразуем вторую сумму справа, обозначив ее через $L$. Используя $s$


раз варьированные формулы Гаусса (4), имеем


\begin{align*}


L&=


\sum_{\alpha=1}^{s}V^{\alpha}\sum_{\beta=1}^{s-\alpha}[\delta^{\beta}


\Phi^k_i\delta^{s-\alpha-\beta}e_k+


\delta^{s-\alpha-\beta}\Phi^k_i\delta^{\beta} e_k+\\


&+\Delta^{\sigma}(\delta^{\beta}\omega^{\sigma}_i


\delta^{s-\alpha-\beta}\nu_{\sigma}+


\delta^{s-\alpha-\beta}\omega^{\sigma}_i


\delta^{\beta}\nu_{\sigma})]=\\


&=\sum_{\alpha=1}^{s-1}\delta^{\alpha}\Phi^k_i


\sum_{\beta=1}^{s-\alpha}V^{\beta}\delta^{s-\alpha-\beta}e_k+\\


&+\sum_{\alpha=2}^{s}\delta^{s-\alpha}\Phi^k_i


\sum_{\beta=1}^{\alpha}V^{\beta}\delta^{\alpha-\beta}e_k-


\sum_{\alpha=2}^{s}\delta^{s-\alpha}\Phi^k_iV^{\alpha}e_k+\\


&+\Delta^{\sigma}\biggl(\,\sum_{\alpha=1}^{s-1}\delta^{\alpha}


\omega^{\sigma}_i\sum_{\beta=1}^{s-\alpha}V^{\beta}


\delta^{s-\alpha-\beta}\nu_{\sigma}+


\sum_{\alpha=2}^{s}\delta^{s-\alpha}\omega^{\sigma}_i


\sum_{\beta=1}^{\alpha}V^{\beta}\delta^{\alpha-\beta}


\nu_{\sigma}-


\sum_{\alpha=2}^{s}\delta^{s-\alpha}\omega^{\sigma}_i


V^{\alpha}\nu_{\sigma}\biggr).


\end{align*}


По теореме 2 получаем


\begin{align*}


L=&\sum_{\alpha=1}^{s}(\delta^{\alpha}\Phi^k_i\delta^{s-\alpha}e_k+


\delta^{s-\alpha}\Phi^k_i\delta^{\alpha}e_k)-\delta^s\Phi^k_ie_k-


\delta^{s-1}\Phi^k_i\delta^1e_k-\\


&-\sum_{\alpha=2}^{s}\delta^{s-\alpha}\Phi^k_iV^{\alpha}e_k+


\Delta^{\sigma}\biggr[\,\sum_{\alpha=1}^{s}(\delta^{\alpha}


\omega^{\sigma}_i\delta^{s-\alpha}\nu_{\sigma}+


\delta^{s-\alpha}\omega^{\sigma}_i\delta^{\alpha}\nu_{\sigma})-\\


&-\delta^s\omega^{\sigma}_i\nu_{\sigma}-


\delta^{s-1}\omega^{\sigma}_i\delta^1\nu_{\sigma}-


\sum_{\alpha=1}^{s}\delta^{s-\alpha}\omega^{\sigma}_iV^{\alpha}


\nu_{\sigma}\biggr].


\end{align*}


%


Вновь применяя $s$ раз варьированные формулы Гаусса, находим


$$


L = \delta^sde_i-\sum_{\alpha=1}^{s-1}(


\delta^{\alpha}\Phi^k_iV^{s-\alpha}e_k+\Delta^{\sigma}


\delta^{\alpha}\omega^{\sigma}_iV^{s-\alpha}\nu_{\sigma})-


(\delta^s\Phi^k_ie_k+\Delta^{\sigma}\delta^s


\omega^{\sigma}_i\nu_{\sigma}).


$$


Подставляя в (21), приходим к тождеству


\begin{equation}


\sum_{\alpha=1}^{s}dV^{\alpha}\delta^{s-\alpha}e_i= 


\sum_{\alpha=1}^{s-1}(\delta^{\alpha}\Phi^k_iV^{s-\alpha}e_k+


\Delta^{\sigma}\delta^{\alpha}\omega^{\sigma}_iV^{s-\alpha}


\nu_{\sigma})+


\delta^s\Phi^k_ie_k+\Delta^{\sigma}\delta^s\omega^{\sigma}_i


\nu_{\sigma},


\end{equation}


где $i,k= 1,\dotsc,n$, $\sigma= 1,\dotsc,p$, $s=1,\dotsc,l$. При


$s= 1$ первая сумма справа исчезает. Отсюда имеем


\begin{equation}


\begin{split}


dV^se_i &=\sum_{\alpha=1}^{s-1}[V^{\alpha}(\delta^{s-\alpha}\Phi^k_ie_k+


\Delta^{\tau}\delta^{s-\alpha}\omega^{\tau}_i\nu_{\tau})-


dV^{\alpha}\delta^{s-\alpha}e_i]+\\


&+\delta^s\Phi^k_ie_k+\Delta^{\tau}\delta^s\omega^{\tau}_i\nu_{\tau}, 


\ i,k=1,\dotsc,n, \ \tau=1,\dotsc,p.


\end{split}


\end{equation}





Дифференцируя второе из равенств (12), будем иметь


$$


d\delta^s\nu_{\sigma}=\sum_{\alpha=1}^{s}dV^{\alpha}\delta^{s-\alpha}


\nu_{\sigma}+\sum_{\alpha=1}^{s}V^{\alpha}d\delta^{s-\alpha}\nu_{\sigma}.


$$


По аналогии с преобразованиями суммы $L$, используя варьированные


формулы Вейнгартена, для второй суммы справа находим


$$


\sum_{\alpha=1}^{s}V^{\alpha}d\delta^{s-\alpha}\nu_{\sigma}=


d\delta^s\nu_{\sigma}+\sum_{\alpha=1}^{s}\Delta^i\delta^{s-\alpha}


\omega^{\sigma}_iV^{\alpha}e_i-


\sum_{\alpha=1}^{s}\Delta^{\tau}\delta^{s-\alpha}\varkappa^{\tau}_{\sigma}


V^{\alpha}\nu_{\tau},


$$


где $i= 1,\dotsc,n$, $\tau,\sigma= 1,\dotsc,p$, $s= 1,\dotsc,l$. Отсюда


и из (23) получаем


\begin{equation}


\begin{split}


dV^s\nu_{\tau}&=


\sum_{\alpha=1}^{s-1}[V^{\alpha}(-\Delta^i\delta^{s-\alpha}


\omega^{\tau}_ie_i+


\Delta^{\rho}\delta^{s-\alpha}\varkappa^{\rho}_{\tau}\nu_{\rho})-


dV^{\alpha}\delta^{s-\alpha}\nu_{\tau}]-\\


&-\Delta^i\delta^s\omega^{\tau}_ie_i+\Delta^{\rho}\delta^s


\varkappa^{\rho}_{\tau}\nu_{\rho}, \ i,k=1,\dotsc,n, 


\ \tau,\rho=1,\dotsc,p.


\end{split}


\end{equation}


Умножая (23) скалярно на $e_i$, $\nu_{\sigma}$, учитывая


(20), получим выражения (19) для коэффициентов


$A^{ij}_s$ и $B^{i\sigma}_s$. Умножая (24) скалярно на


$\nu_{\sigma}$, получим требуемое выражение для коэффициентов


$C^{\tau\sigma}_s$. Равенство (17) получится, если в этих


рассуждениях положить $s= 1$.


\end{pf}





\begin{thebibliography}{99}





\bibitem{Efi}


Ефимов Н.В. {\em Качественные вопросы теории деформаций поверхностей} //


УМН. -- 1948. -- Т.\,3. -- \No\,2. -- С.\,47--158.





\bibitem{Sab} Сабитов И.Х.


{\em Локальная теория изгибаний поверхностей} // Итоги науки и техн.


ВИНИТИ. Современ. пробл. матем. -- 


1989. -- Т.\,48. -- С.\,196--270.





\bibitem{Iva1} Иванова-Каратопраклиева И., Сабитов И.Х.


{\em Изгибание поверхностей}. I // Итоги науки и техн. ВИНИТИ. 


Пробл. геометрии. -- 1991. -- T.\,23. -- С.\,131--184.





\bibitem{Iva2} Иванова-Каратопраклиева И., Сабитов И.Х.


{\em Изгибание поверхностей}. II // Итоги науки и техн. ВИНИТИ. 


Современ. матем. и ее прилож. -- 1995. -- T.\,8. -- С.\,108--167.





\bibitem{Car1} Cartan E. {\em Sur les d\'eformations des


hypersurfaces dans l'espace conforme de dimension $\ge5$} // Bull. Soc.


math. France. -- 1917. V.\,45. -- P.\,57--121.





\bibitem{DoC} Do Carmo M, Dajczer M. {\em Conformal


Rigidity} // Amer. J. Math. -- 1987. -- V.\,109. -- P.\,963--985.





\bibitem{Fom} Фоменко В.Т., Бикчантаев И.А.


{\em Применение обобщенных аналитических функций на римановых поверхностях


к исследованию $G$-деформаций двумерных поверхностей в $E^4$}


// Матем. сб. -- 1988. -- Т.\,136. -- \No\,4. -- C.\,561--573.





\bibitem{Bes} Бескоровайная Л.Л. {\em О  бесконечно


малых ареальных деформациях овальных поверхностей}


// Изв. вузов. Математика. -- 1983. -- \No\,5. -- С.\,69--71.





\bibitem{Yano} Yano K. {\em Sur la th\'eorie des d\'eformations


infinit\'esimales} // J. Fac. Sci. Univ. Tokyo. Ser.\,I -- 1949.


-- \No\,6. -- P.\,1--75.





\bibitem{Chen} Chen  B.Y., Yano K. {\em On the theory of


normal variations}


// J. Different. Geom. -- 1978. -- V.\,13. -- \No\,1. -- P.\,1--10.





\bibitem{Dar} Darboux G. {\em Le\c{c}ons sur la th\'eorie des


surfaces et les applications g\'eometriques du calcul infinit\'esimal}.


4 partie. -- Paris.: Gauthier-Villars, Imprimeur-\'editeur, 1946. -- 554 p.





\bibitem{Mark1} Марков П.Е. {\em Бесконечно малые изгибания


одного класса многомерных поверхностей с краем} // Матем. сб.


-- 1983. -- Т.\,121. -- \No\,1. -- С.\,48--59.





\bibitem{Mark2} Марков П.Е. {\em Бесконечно малые изгибания


высших порядков многомерных поверхностей в пространствах постоянной


кривизны} // Матем. сб. 1987. -- Т.\,133. -- \No\,1. -- С.\,64--85.





\bibitem{Vek} Векуа И.Н. {\em Обобщенные аналитические функции}. -- 


М.: 1988. -- 510 с.





\bibitem{Mark3} Марков П.Е. {\em Бесконечно малые изгибания


некоторых многомерных поверхностей}


// Матем. заметки. -- 1980. -- Т.\,27. -- \No\,3. -- 


С.\,469--479.





\bibitem{Mark4} Марков П.Е. {\em О погружении метрик, близких к


погружаемым} // Укр. геом. сб. -- 1992. -- \No\,35. -- С.\,49--67.





\bibitem{All} Allendoerfer C.B. {Rigidity for spaces of


class greater than one} // Amer. J. Math. -- 1939. -- \No\,61.


-- P.\,633--644.





\bibitem{Yan} Яненко Н.Н. {\em Некоторые необходимые признаки


изгибаемых поверхностей $V_m$ в $(m+q)$-мерном евклидовом пространстве}


// Тр. семин. по векторн. и тензорн. анализу. -- 1952. -- \No\,9.


-- C.\,326--287.





\bibitem{Kow} Kowalski O. {\em Some algebraic theorems on


vector-valued forms and their geometric applications} // Colloq. Math.


-- 1972. -- T.\,26. -- S.\,59--92.





\bibitem{Gard} Gardner R.B. {\em New viewpoints in the


geometry of submanifolds of $R^N$} // Bull. Amer. Math. Soc.


-- 1977. -- V.\,83. -- \No\,1. -- P.\,1--35.





\bibitem{Chern} Chern S.-S., Osserman R. {\em Remarks on


the Riemannian metric of a minimal submanifold} // Lect. Notes Math.


-- 1981. -- \No\,894. -- P.\,49--90.





\bibitem{Kob} Кобаяси Ш., Номидзу К. {\em Основы


дифференциальной геометрии}. Т 2. -- 


М.: Наука, 1981. -- 414 с.





\bibitem{Vish} Вишневский В.В.,


Широков А.П., Шурыгин В.В. {\em Пространства над алгебрами}.


-- Казань: Изд-во Казанск. ун-та, 1985. -- 262 с.





\bibitem{Zu} Зуланке Р., Винтген П.


{\em Дифференциальная геометрия и расслоения}. --


М.: Мир, 1975. -- 352~с.





\bibitem{Weyl} Вейль Г. {\em Классические группы. Их инварианты


и представления}. -- М.: ГИИЛ, 1947. -- 408~с.





\bibitem{Gan} Гантмахер Ф.P. {\em Теория матриц}. -- 3-е изд. --


М.: Наука, 1967. -- 575 с.





\bibitem{Post} Постников М.М. {\em Группы и алгебры Ли}.


Учеб. пособие. -- М.: Наука, 1982. -- 447 с.





\bibitem{Wolf} Вольф Дж. {\em Пространства постоянной кривизны}. --


М.: Наука, 1982. -- 480 с.





\bibitem{Eis} Эйзенхарт Л.П. {\em Риманова геометрия}. --


М.: ГИИЛ, 1948. -- 316 с.





\bibitem{DeRh} Де Рам Ж. {\em Дифференцируемые многообразия}. -- 


М.: ГИИЛ, 1956. -- 250 с.





\bibitem{Bor1} Боровский Ю.Е. {\em Вполне  интегрируемые


системы Пфаффа} // Изв. вузов. Математика. -- 1959. -- \No\,3. -- С.\,28--40.





\bibitem{Bor2} Боровский Ю.Е. {\em О вполне интегрируемых


системах Пфаффа} // Изв. вузов.


Математика. -- 1960. -- \No\,1. -- С.\,35--38.





\bibitem{Bor3} Боровский  Ю.Е.  {\em Системы  Пфаффа  с


коэффициентами из $L_n$ и их геометрические приложения} // Сиб. матем.


журн. -- 1988. -- Т.29. -- \No\,2. -- С.\,10--16.





\bibitem{Mark5} Марков П.Е. {\em Бесконечно малые изгибания


высших порядков многомерных поверхностей} // Укр. геом. сб. -- 1982.


-- \No\,25. -- С.\,87--94.





\bibitem{Kli} Климентов С.Б. {\em О продолжении бесконечно


малых изгибаний высших порядков односвязной поверхности положительной


кривизны} // Матем. заметки. -- 1984. -- Т.\,36. -- \No\,3. -- C.\,393--403.





\bibitem{Car2} Картан Э. {\em Геометрия римановых пространств}. --


М.: ОНТИ НКТП, 1936. -- 244 с.





\bibitem{Per} Перепелкин Д.И. {\em Кривизна и нормальные


пространства многообразия $V_m$ в $R_n$} // Матем. сб. -- 1935.


-- Т.\,42. -- \No\,1--3. -- C.\,100--120.





\end{thebibliography} 


 


\vskip 2cm 


 


\post{\it Ростовский государственный университет}{\it Поступила}


\post{}{27.03.1995} 


 


\end{document} 


